\section{The load is geometric at every time}
\label{dist:sec:geometric}

The generating function \cref{dist:eq:gcat} is a ratio of two linear
functions of $z$, and that fact alone decides the shape of the whole
distribution. Extracting the coefficients returns not an unstructured sequence
but a geometric law, at every time, with a ratio that moves. Naming that ratio
is the single most useful thing this chapter does, because every result after
this section is a statement about where it goes.

Coefficients come out of a PGF by repeated differentiation at the origin,
$p_n=\frac{1}{n!}\partial_z^nG|_{z=0}$; the construction is given in
\mextref{m:app:extract}{the coefficient-extraction appendix of \ChM} at
\mextref{m:eq:stateProb}{the state-probability formula there}. Applied to
\cref{dist:eq:gcat}, the derivatives close on themselves,
\begin{equation}
\pdd{^n G}{z^n}
= n!\,(a-b)^2 w\,
\frac{(1-w)^{n-1}}{\bigl(b-aw-(1-w)z\bigr)^{n+1}},
\label{dist:eq:dnG}
\end{equation}
so that setting $z=0$ and dividing by $n!$ gives, for $n\ge1$,
\begin{equation}
p_n(t)=\lrb{\frac{a-b}{b-aw}}^{2}w
\lrb{\frac{1-w}{b-aw}}^{n-1}.
\label{dist:eq:stateprob}
\end{equation}
The empty state is the only one that sees the convention. Under catastrophe
$p_0=G(0,t)=ab(1-w)/(b-aw)$, the probability of reaching $0$ by the death of
every individual; under killing it carries the extra $1-I(t)$ of
\cref{dist:prop:conventions}, and nothing else changes.

\subsection{The sliding ratio}
\label{dist:sec:sliding}

There is more structure in \cref{dist:eq:stateprob} than it first shows. The
entire dependence on $n$ sits in one factor. Write
\begin{equation}
P(t):=\frac{1-w}{b-aw},
\qquad w=\ee^{\theta t},
\label{dist:eq:pdef}
\end{equation}
and the state probabilities become $p_n(t)=p_1(t)P(t)^{n-1}$.

\begin{theorem}[The load is geometric at every time]
\label{dist:thm:geometric}
For every $t>0$, conditional on the load being non-zero, $\Xt$ is
geometrically distributed on $\{1,2,\ldots\}$ with ratio $P(t)$:
\begin{equation}
p_n(t)=p_1(t)\,P(t)^{n-1},\qquad n\ge1.
\label{dist:eq:geometric}
\end{equation}
The ratio increases monotonically in $t$, from
\begin{equation}
P(0)=0 \qquad\text{to}\qquad P(\infty)=\frac1a .
\label{dist:eq:plimits}
\end{equation}
\end{theorem}

\begin{proof}
Equation \cref{dist:eq:geometric} is \cref{dist:eq:stateprob} read with the
$n$-dependent factor named. For the limits, $w=1$ at $t=0$ gives $P(0)=0$,
and $w\to\infty$ gives $P\to1/a$; monotonicity follows from
$P'(t)=\lambda(a-b)^2w/(b-aw)^2>0$, since $b-aw$ never vanishes for $t\ge0$.
\end{proof}

What $P$ is, is worth saying plainly. It is the price in probability of one
more particle: $p_{n+1}/p_n=P(t)$ for every $n\ge1$, so each additional unit
of load is the same factor less likely than the one below it, whatever that
load happens to be. At founding $P=0$ and the cell holds its single particle
with certainty. Thereafter $P$ rises, the geometric flattens, and the mean
load conditional on non-fixation, $1/(1-P)$, rises with it --- from $1$ at
founding towards $a/(a-1)$. That the rise never reverses is what
\cref{dist:thm:geometric} establishes, and it is what makes the endpoint well
defined: a cell that has survived longer is spread over more states, never
fewer.

So the transient laws form a one-parameter family of geometrics sliding
towards a fixed limit, and $P$ is the only thing in them that moves.
\Cref{dist:fig:slide} shows the slide: on a logarithmic scale each law is a
straight line, and the lines rotate onto a limiting slope rather than
deforming. The limit $1/a$ is where the next three sections live ---
\cref{dist:sec:qsd} takes it, \cref{dist:sec:burst-size} integrates against
it, \cref{dist:sec:burst-time} conditions on it.
\Cref{dist:fig:readings} is that sentence drawn.

\begin{figure}[tbp]
\centering
\includegraphics[width=0.91\textwidth]{fig02_three_readings/fig02}
\caption{One endpoint, three readings. The ratio $P(t)$ of
\cref{dist:thm:geometric} slides from $0$ at founding to $1/a$, and each of
the next three sections applies a different operation to that endpoint:
\cref{dist:sec:qsd} takes it as the quasi-stationary ratio,
\cref{dist:sec:burst-size} integrates the rupture hazard against it, and
\cref{dist:sec:burst-time} conditions the rupture time on the size it
produces. The first two return the same law, differing only by the
eventual-burst probability $1-b$; that coincidence is
\cref{dist:thm:identity}.}
\label{dist:fig:readings}
\end{figure}

\begin{figure}[tbp]
\centering
\includegraphics[width=0.82\textwidth]{fig03_geometric_slide/fig03}
\caption{State probabilities $p_n(t)$ at successive times for
$(\lambda,\mu,\delta)=(1,0.2,0.05)$. On the logarithmic vertical scale each
geometric law is a straight line whose ratio $P(t)$ increases towards $1/a$;
the inset shows that approach directly. The curves fall in $n$ and rise in
$t$: the load spreads as $P$ slides.}
\label{dist:fig:slide}
\end{figure}

\subsection{Normalisation: the sum is \texorpdfstring{$\Ifix$}{Ifix}}
\label{dist:sec:sum-ihat}

One consequence is immediate, and it checks the closed forms of
\cref{dist:sec:recap} against a derivation that shares nothing with them.

\begin{corollary}[The positive states carry the non-fixation probability]
\label{dist:cor:sum}
For every $t\ge0$,
\begin{equation}
\sum_{n\ge1}p_n(t)
=\frac{(a-b)^2 w}{\lrb{B+Aw}\lrb{aw-b}}
=\Ifix(t).
\label{dist:eq:pnsum}
\end{equation}
\end{corollary}

\begin{proof}
By \cref{dist:thm:geometric} the sum is the geometric series
$p_1(t)/(1-P(t))$. Now
\[
1-P(t)=\frac{b-aw-1+w}{b-aw}
=\frac{-(1-b)-(a-1)w}{b-aw}
=\frac{B+Aw}{aw-b},
\]
and dividing $p_1(t)=(a-b)^2w/(b-aw)^2$ by this gives the middle expression,
which is the closed form for $\Ifix$ in \cref{dist:eq:idihat}.
\end{proof}

The sum of the positive state probabilities is exactly the probability of
being neither ruptured nor extinct. \ChCorecap{} reached the same $\Ifix$ by
decomposing absorption rather than by summing a distribution; that two such
different routes agree is a check on both.
